\documentclass[10pt,a4paper]{amsart}
\usepackage[margin=19mm]{geometry}
\usepackage{amsmath,amssymb,mathtools}
\usepackage[scr=rsfs,cal=euler]{mathalfa}
\usepackage{xfrac}
\usepackage[unicode,hidelinks,bookmarks=false]{hyperref}
\newcommand{\ie}{i.e.\ }
\newcommand{\cf}{cf.\ }
\newcommand{\resp}{respectively\ }
\newcommand{\Subsection}{\subsection}
\providecommand{\nobreakdash}{\nobreak\hbox{-}}
\setlength{\emergencystretch}{2em}
\multlinegap0pt
\usepackage[shortlabels]{enumitem}
\usepackage{amssymb}
\usepackage{mathtools}
\usepackage{bm}
\newtheorem{theorem}{Theorem}
\newtheorem{lemma}{Lemma}
\usepackage{cleveref}
\crefname{theorem}{Theorem}{Theorems}
\crefname{lemma}{Lemma}{Lemmas}
\newcommand{\Pcal}{\mathbb{P}}
\newcommand{\R}{\mathbb{R}}
\newcommand{\abs}[1]{\left\lvert #1\right\rvert}
\newcommand{\dd}{\,\mathrm{d}}
\newcommand{\diver}{\operatorname{div}}
\newcommand{\ip}[2]{\left\langle #1,#2\right\rangle}
\newcommand{\norm}[1]{\left\lVert #1\right\rVert}
\title{Nonlocal Hyperdissipative Perturbations of the Three-Dimensional Navier--Stokes System}
\author{Veli Shahmurov, Rishad Shahmurov}
\date{Source: arXiv:2604.04167v1 (CC BY 4.0)}
\begin{document}
\maketitle

\noindent\emph{Source and licence.} Nonlocal Hyperdissipative Perturbations of the Three-Dimensional Navier--Stokes System. Version arXiv:2604.04167v1 (CC BY 4.0), distributed by arXiv under the Creative Commons Attribution 4.0 International licence, http://creativecommons.org/licenses/by/4.0/, as recorded in the arXiv licence metadata for this exact version and shown on the abstract page https://arxiv.org/abs/2604.04167v1. Full licence notice: \texttt{licenses/arxiv-2604.04167-CC-BY-4.0.txt}.\par\smallskip

\noindent\textit{Editorial scope.} Counted unit: the article's theorem thm:lions (source0.tex lines 746--748). The article's complete original proof of it is delivered with it (source0.tex lines 750--845, one \texttt{proof} environment of 96 lines). No mathematics is written, repaired or re-derived: the statement and the proof are the article's own lines, in the article's own notation, and no retained line is substituted or re-flowed.  Retained uncounted as the source-local context the proof needs: lines 149--159; lines 180--183; lines 188--196; lines 240--258; lines 618--624; lines 709--720.  Nothing was omitted from any retained span, so the package carries no omission record.  No retained line contains a citation, so the wrapper carries no bibliography and no bibliography entry.  No retained line contains a figure environment, an image-inclusion command or drawing code, so no image is delivered and the package declares no asset dependency.  Editorial formatting only, in addition: an AMS \texttt{amsart} preamble that restates the article's own declarations and macros used by the retained lines, namely the environments \texttt{theorem}, \texttt{lemma} on separate counters, together with the standard packages those lines need.  The retained environments print in this document as Theorem 1, Lemma 1 in order of appearance, whereas the article prints the same items with the numbering of its own full document.  The complete native source of the article as distributed in the arXiv e-print for this exact version is bundled unchanged as \texttt{sources/197866/source0.tex}.
\begin{equation}\label{eq:main}
\left\{
\begin{aligned}
&\partial_t u + \Pcal \diver (u\otimes u) + \nu(-\Delta)u + M u = 0,
\\
&\diver u = 0,
\\
&u|_{t=0}=u_0,
\end{aligned}
\right.
\end{equation}

\begin{equation}\label{eq:symbol-growth}
c_0 |\xi|^{2\alpha} \le m(\xi) \le c_1 (1+|\xi|^{2\alpha}),
\qquad \xi\in\R^3.
\end{equation}

\begin{equation}\label{eq:coercive-equivalence}
c_0 \norm{\Lambda^\alpha u}_{L^2}^2
\le
\ip{M u}{u}_{L^2}
=
\int_{\R^3} m(\xi)\abs{\hat u(\xi)}^2 \dd \xi
\le
c_1 \bigl(\norm{u}_{L^2}^2 + \norm{\Lambda^\alpha u}_{L^2}^2\bigr).
\end{equation}

\begin{theorem}[Main results]\label{thm:main}
Assume that $m\in\mathfrak M_\alpha$ for some $\alpha>1$.
\begin{enumerate}[label=(\roman*)]
\item For smooth solutions the exact energy identity
\[
\frac12 \frac{\dd}{\dd t}\norm{u(t)}_{L^2}^2
+
\nu \norm{\nabla u(t)}_{L^2}^2
+
\int_{\R^3} m(\xi)\abs{\hat u(\xi,t)}^2 \dd \xi
=0
\]
holds for every $t>0$.
\item For every $u_0\in L^2_\sigma(\R^3)$ there exists a global weak solution of \eqref{eq:main}.
\item If $s>\frac52$ and $u_0\in H^s_\sigma(\R^3)$, then there exists $T=T(\norm{u_0}_{H^s})>0$ and a unique strong solution on $[0,T]$.
\item If $\alpha\ge \frac54$ and $u_0\in H^s_\sigma(\R^3)$ with $s>\frac52$, then the strong solution is global.
\item If $\alpha>1$ and $u_0\in H^s_\sigma(\R^3)$ with $s>\frac52$ is sufficiently small, then the strong solution is global.
\end{enumerate}
\end{theorem}

\begin{theorem}[Local strong well-posedness]\label{thm:local}
Assume that $m\in\mathfrak M_\alpha$ with $\alpha>1$, let $s>\frac52$, and let $u_0\in H^s_\sigma(\R^3)$. Then there exists $T=T(\norm{u_0}_{H^s})>0$ and a unique strong solution
\[
u\in C([0,T];H^s_\sigma(\R^3))\cap L^2(0,T;H^{s+\alpha}_\sigma(\R^3))
\]
of \eqref{eq:main}. Moreover, the solution depends continuously on $u_0$ in $H^s$.
\end{theorem}

\begin{lemma}[Critical product estimate]\label{lem:critical-product}
Let $s>\frac52$ and $\alpha\ge \frac54$. Then there exists $C=C(s,\alpha)$ such that for every pair of divergence-free vector fields $u,v\in H^s_\sigma(\R^3)$,
\begin{equation}\label{eq:critical-product}
\norm{\Pcal\diver(u\otimes v)}_{H^{s-\alpha}}
\le C\norm{u}_{H^\alpha}\norm{v}_{H^s} + C\norm{u}_{H^s}\norm{v}_{H^\alpha}.
\end{equation}
In particular,
\begin{equation}\label{eq:critical-product-diagonal}
\norm{\Pcal\diver(u\otimes u)}_{H^{s-\alpha}}
\le C\norm{u}_{H^\alpha}\norm{u}_{H^s}.
\end{equation}
\end{lemma}

\begin{theorem}[Global strong solutions at and above the Lions exponent]\label{thm:lions}
Assume that $m\in\mathfrak M_\alpha$ with $\alpha\ge \frac54$ and let $s>\frac52$. Then every initial datum $u_0\in H^s_\sigma(\R^3)$ generates a unique global strong solution of \eqref{eq:main}.
\end{theorem}

\begin{proof}
Let $u$ denote the unique strong solution supplied by \cref{thm:local} on its maximal interval of existence $[0,T_*)$, where $T_*\in(0,\infty]$. We shall prove that $T_*=\infty$.

\smallskip
\noindent\textbf{Step 1: high-order energy identity.}
Apply $\Lambda^s$ to \eqref{eq:main} and take the $L^2$ inner product with $\Lambda^s u$. Since $\Pcal$ is selfadjoint on divergence-free fields and commutes with Fourier multipliers,
\begin{equation}\label{eq:Hs-energy-start}
\frac12\frac{\dd}{\dd t}\norm{u}_{H^s}^2
+\nu\norm{\nabla u}_{H^s}^2
+\int_{\R^3}(1+|\xi|^2)^s m(\xi)|\hat u(\xi)|^2\,\dd\xi
=-\ip{\Lambda^s \Pcal\diver(u\otimes u)}{\Lambda^s u}.
\end{equation}
Using the lower bound in \eqref{eq:symbol-growth}, we obtain
\begin{align}
\int_{\R^3}(1+|\xi|^2)^s m(\xi)|\hat u(\xi)|^2\,\dd\xi
&\ge c_0\int_{\R^3}(1+|\xi|^2)^s|\xi|^{2\alpha}|\hat u(\xi)|^2\,\dd\xi \notag\\
&\ge c\norm{u}_{H^{s+\alpha}}^2-C\norm{u}_{H^s}^2,
\label{eq:Hs-diss-lower}
\end{align}
where the last step is the elementary low-frequency comparison
\[
(1+|\xi|^2)^s|\xi|^{2\alpha}\ge c(1+|\xi|^2)^{s+\alpha}-C(1+|\xi|^2)^s.
\]

\smallskip
\noindent\textbf{Step 2: estimate of the nonlinear term.}
Rewrite the right-hand side of \eqref{eq:Hs-energy-start} by shifting $\alpha$ derivatives:
\[
\bigl|\ip{\Lambda^s \Pcal\diver(u\otimes u)}{\Lambda^s u}\bigr|
=
\bigl|\ip{\Lambda^{s-\alpha}\Pcal\diver(u\otimes u)}{\Lambda^{s+\alpha}u}\bigr|.
\]
Therefore, by \cref{lem:critical-product},
\begin{align}
\bigl|\ip{\Lambda^s \Pcal\diver(u\otimes u)}{\Lambda^s u}\bigr|
&\le C\norm{\Pcal\diver(u\otimes u)}_{H^{s-\alpha}}\norm{u}_{H^{s+\alpha}} \notag\\
&\le C\norm{u}_{H^\alpha}\norm{u}_{H^s}\norm{u}_{H^{s+\alpha}} \notag\\
&\le \frac{c}{2}\norm{u}_{H^{s+\alpha}}^2 + C\norm{u}_{H^\alpha}^2\norm{u}_{H^s}^2.
\label{eq:nonlinear-Hs-est}
\end{align}
Combining \eqref{eq:Hs-energy-start}, \eqref{eq:Hs-diss-lower}, and \eqref{eq:nonlinear-Hs-est}, and absorbing the $\frac c2\norm{u}_{H^{s+\alpha}}^2$ term into the left-hand side, we arrive at
\begin{equation}\label{eq:Hs-differential}
\frac{\dd}{\dd t}\norm{u}_{H^s}^2 + c\norm{u}_{H^{s+\alpha}}^2
\le C\bigl(1+\norm{u}_{H^\alpha}^2\bigr)\norm{u}_{H^s}^2.
\end{equation}

\smallskip
\noindent\textbf{Step 3: control of the coefficient by the basic energy law.}
Since the strong solution satisfies the exact $L^2$ energy identity from \cref{thm:main}(i),
\begin{equation}\label{eq:lions-basic-energy}
\sup_{0\le t<T}\norm{u(t)}_{L^2}^2
+2\nu\int_0^T\norm{\nabla u(t)}_{L^2}^2\,\dd t
+2\int_0^T\ip{Mu(t)}{u(t)}\,\dd t
\le \norm{u_0}_{L^2}^2
\end{equation}
for every $T<T_*$. By \eqref{eq:coercive-equivalence},
\[
\int_0^T \norm{\Lambda^\alpha u(t)}_{L^2}^2\,\dd t
\le C\norm{u_0}_{L^2}^2.
\]
Consequently,
\begin{equation}\label{eq:alpha-time-integrable}
\int_0^T \norm{u(t)}_{H^\alpha}^2\,\dd t
\le C\Bigl(T\norm{u_0}_{L^2}^2+\norm{u_0}_{L^2}^2\Bigr)
<\infty
\qquad\text{for every }T<T_*.
\end{equation}
Thus the coefficient on the right-hand side of \eqref{eq:Hs-differential} is integrable on every finite subinterval of $[0,T_*)$.

\smallskip
\noindent\textbf{Step 4: Gronwall bound and continuation.}
Fix $T<T_*$. Integrating \eqref{eq:Hs-differential} and applying Gronwall's inequality gives
\begin{equation}\label{eq:Hs-global-bound}
\sup_{0\le t\le T}\norm{u(t)}_{H^s}^2
\le
\norm{u_0}_{H^s}^2
\exp\Bigl(C T + C\int_0^T\norm{u(t)}_{H^\alpha}^2\,\dd t\Bigr).
\end{equation}
Using \eqref{eq:alpha-time-integrable}, the right-hand side is finite and depends only on $T$, $\norm{u_0}_{L^2}$, and $\norm{u_0}_{H^s}$. In particular,
\[
\sup_{0\le t<T}\norm{u(t)}_{H^s}<\infty
\qquad	ext{for every finite }T<T_*.
\]
Moreover, integrating \eqref{eq:Hs-differential} once more yields
\[
\int_0^T\norm{u(t)}_{H^{s+\alpha}}^2\,\dd t<\infty.
\]

Suppose for contradiction that $T_*<\infty$. Choose $t_0\in(0,T_*)$ close enough to $T_*$ that
\[
\norm{u(t_0)}_{H^s}\le 2\sup_{0\le t<T_*}\norm{u(t)}_{H^s}<\infty.
\]
By \cref{thm:local}, the lifespan of a strong solution launched from the datum $u(t_0)$ depends only on its $H^s$ norm. Hence there exists $\tau>0$, depending only on the preceding uniform bound, such that the solution starting from $u(t_0)$ extends uniquely to $[t_0,t_0+\tau]$. Taking $t_0$ so close to $T_*$ that $t_0+\tau>T_*$ contradicts the maximality of $T_*$. Therefore $T_*=\infty$.

Uniqueness is already part of the local theory and therefore carries over to the global solution.
\end{proof}

\end{document}
